\section{Generator Inversion}
\label{section:inversion}

\subsection{Algorithm}
We present an implementation of Strassen's inversion method
specifically adapted for displacement generators.
Strassen's approach relies on recursive block-matrix partitioning and
Schur complements,

\begin{algorithm}[H]
  \caption{Strassen-type Inversion for Displacement Generators}
  \begin{algorithmic}[1]
    \Procedure{StrassenInverseGenerators}{$G_A, H_A$}

    \caption{Strassen-type Inversion for Displacement Generators
    (adapted from \cite{aecf-2017-pdf})}

    \If{$n = 1$} \Comment{\textbf{Step 1: Base Case}}

    \State \textbf{return} $G_D$ $H_D$ representing the inverse of
    the scalar value of $A$ from $G_A, H_A$

    \EndIf

    \Statex

    \State $G_a, H_a, \dots, G_d, H_d \gets
    \textsc{SplitQuadrants}(G_A, H_A)$ \Comment{\textbf{Step 2: Split
    into Quadrants}}

    \Statex \Comment{\textbf{Step 3: Recursion\#1}}

    \State $G_e, H_e \gets \textsc{StrassenInverseGenerators}(G_a,
    H_a)$ \Comment{$e := a^{-1}$}

    \Statex

    \State $G_{ce}, H_{ce} \gets \textsc{Compress}(\textsc{Mul}(G_c,
    H_c, G_e, H_e))$ \Comment{\textbf{Step 4: Schur Complement}}

    \State $G_{eb}, H_{eb} \gets \textsc{Compress}(\textsc{Mul}(G_e,
    H_e, G_b, H_b))$

    \State $G_{ceb}, H_{ceb} \gets
    \textsc{Compress}(\textsc{Mul}(G_c, H_c, G_{eb}, H_{eb}) \times -1)$

    \State $G_S, H_S \gets \textsc{Compress}(\textsc{Add}(G_d, H_d,
    G_{ceb}, H_{ceb}))$ \Comment{$S := d - ceb$}

    \Statex \Comment{\textbf{Step 5: Recursion\#2}}

    \State $G_t, H_t \gets \textsc{StrassenInverseGenerators}(G_S,
    H_S)$ \Comment{$t := S^{-1}$}

    \Statex

    \State $G_{ebt}, H_{ebt} \gets
    \textsc{Compress}(\textsc{Mul}(G_{eb}, H_{eb}, G_t, H_t))$
    \Comment{\textbf{Step 6: Strassen Assembly}}

    \State $G_{tce}, H_{tce} \gets
    \textsc{Compress}(\textsc{Mul}(G_t, H_t, G_{ce}, H_{ce}))$

    \State $G_{ebtce}, H_{ebtce} \gets
    \textsc{Compress}(\textsc{Mul}(G_{ebt}, H_{ebt}, G_{ce}, H_{ce}))$

    \Statex

    \State $G_x, H_x \gets \textsc{Compress}(\textsc{Add}(G_e, H_e,
    G_{ebtce}, H_{ebtce}))$ \Comment{$x := e + ebtce$}

    \State $G_y, H_y \gets \textsc{Negate}(G_{ebt}, H_{ebt})$
    \Comment{$y := -ebt$}

    \State $G_z, H_z \gets \textsc{Negate}(G_{tce}, H_{tce})$
    \Comment{$z := -tce$}

    \Statex

    \State $G_{Inv}, H_{Inv} \gets \textsc{PackQuadrants}(G_x, H_x,
    G_y, H_y, G_z, H_z, G_t, H_t)$ \Comment{\textbf{Step 7: Pack}}

    \State \textbf{return} \textsc{Compress}($G_{Inv}, H_{Inv}$)

    \EndProcedure
  \end{algorithmic}
\end{algorithm}

This algorithm is based on Algorithm 10.1 from \cite[Algorithme
10.1]{aecf-2017-pdf}.
Notice that we only operate on $\Phi_+$ displacement and bypassing
$\Phi_-$ generators.

\begin{proposition}
  The Strassen-type inversion algorithm maintains $\Phi_+$
  displacement generators throughout its recursive execution, without
  necessitating conversion to $\Phi_-$ generators.
\end{proposition}

\begin{proof}
  We proceed by induction on the size of the matrix $n = 2k$. Let
  $\Phi_+(X) = X - ZXZ^T$ be the displacement operator.

  \paragraph{Base Case ($n=1$):}
  For a scalar matrix $A = [a]$, the operator $Z$ is a $1 \times 1$ zero matrix.
  The displacement of $A$ is $\Phi_+(A) = A - 0 \cdot A \cdot 0 = A$.
  \paragraph{Inductive Step:}
  Consider a matrix $A$ of size $2k \times 2k$ represented by
  $\Phi_+$ generators.
  Inversion relies on splitting the sub-blocks $a, b, c, d$ which are
  represented by $\Phi_+$ displacements,
  recursive calls to invert top-left block $a$ ($e = a^{-1}$) and
  antoher for the Schur complement $t = S^{-1}$ and each outputs a
  $\Phi_+$ generator.
  Multiplication and addition algorithms are taking and returning a
  generator of $\Phi_+$ displacement.
  Each computation and the final assembly remain in $\Phi_+$ displacement.

  Because the base case returns the inverse to $\Phi_+$ generators,
  and the arithmetics are on $\Phi_+$ generator operations, by
  induction, we maintain $\Phi_+$ displacement during the entire execution.
\end{proof}

\subsection{Split and Pack quadrants from displacement generators}

Any divide and conquer algorithm will divide the problem to create
sub-problems to economise in complexity.
This algorithm calls for a dense matrix $A$ be split into four quadrants:
$$ A =
\begin{pmatrix} a & b \\ c & d
\end{pmatrix} $$

Since we don't have direct access to the dense matrix $A$,
any splitting and packing operation done will be on the generators
$G_A$ and $H_A$ which represent the displacement matrix.
When we attempt to evaluate the local displacement of a sub-block
(for example, $\Phi_+(d)$), we cannot simply extract the rows from
the global generators $G_A$ and $H_A$.
The global generators are corrupted by the boundary elements of $a$
that shifted into $b$ and $c$.
We must correct these boundary crossings which we call bleeding.

\vspace{0.2cm}

\subsubsection{Splitting}

During the split phase, we extract the local generators for
sub-blocks $a, b, c, d$.
We begin by splitting the global generators $G_A, H_A \in
\mathbb{R}^{n \times r}$ into top and bottom blocks of sizes $n_1 =
ceil(n/2)$ and $n_2 = floor(n/2)$:
\begin{equation*}
  G =
  \begin{pmatrix} G_{\text{TOP}} \\ G_{\text{BOTTOM}}
  \end{pmatrix}
  \quad
  H =
  \begin{pmatrix} H_{\text{TOP}} \\ H_{\text{BOTTOM}}
  \end{pmatrix}
\end{equation*}

\vspace{0.5em}
\noindent\textbf{Block $\boldsymbol{a}$:}
This block receives no shifted data from other quadrants.
Its generators are the global top halves:
$$ G_{a} = G_{\text{TOP}}, \quad H_{a} = H_{\text{BOTTOM}} $$

\vspace{0.5em}
\noindent\textbf{Block $\boldsymbol{b}$:}
The last column of block $a$, denoted $v_a$, shifts right into the
first column of block $b$.
To correct this bleed, we must reconstruct $v_a$ from the global generators.
Naively evaluating the displacement sum for a boundary vector
requires $O(n^2 r)$ operations.
However, this evaluation is equivalent to computing the coefficients
of a polynomial multiplication.
By packing the columns of $G$ and the reversed columns of $H$ into
polynomials, we extract $v_a$ with polynomial multiplication in $O(r
\cdot M(n))$ time.
Once extracted, we shift $v_a$ down using the operator $Z_{n_1}$ and
append it to $G_{\text{TOP}}$.
To apply this correction only to the first column of $b$, we append
the vector $e_0^{(n_2)} = (1, 0, \dots, 0)^T$ to $H_{\text{bot}}$:
$$ G_{b} =
\begin{pmatrix} G_{\text{TOP}} & Z_{n_1} v_a
\end{pmatrix}, \quad H_{b} =
\begin{pmatrix} H_{\text{BOTTOM}} & e_0^{(n_2)}
\end{pmatrix} $$

\vspace{0.5em}
\noindent\textbf{Block $\boldsymbol{c}$:}
Bleeding occurs to $c$ from above via the last row of $a$, denoted
$r_a$. Using the same polynomial operation, we recover $r_a$.
Because $H$ generators represent the transposed displacement
components, we apply the shift $Z_{n_1}$ to $r_a$ and append it to
$H_{\text{TOP}}$:
$$ G_{c} =
\begin{pmatrix} G_{\text{BOTTOM}} & e_0^{(n_2)}
\end{pmatrix}, \quad H_{c} =
\begin{pmatrix} H_{\text{TOP}} & Z_{n_1} r_a
\end{pmatrix} $$

\vspace{0.5em}
\noindent\textbf{Block $\boldsymbol{d}$:}
Block $d$ is exposed to bleeding from all directions: the last column
of $c$ ($v_c$) shifting right, the last row of $b$ ($r_b$) shifting down,
and the exact bottom-right scalar of $a$ ($s_a$) shifting diagonally down-right.
We extract $v_c$ and $r_b$ using polynomial operation.
Because $s_a$ is a single cell, polynomial overhead is unnecessary;
it is computed directly in $\O(n r)$ time.
Scaling the vector $e_0^{(n_2)}$ by $s_a$ and addint to the shifted vectors:
$$ G_{d} =
\begin{pmatrix} G_{\text{BOTTOM}} & s_a e_0^{(n_2)} + Z_{n_2} v_c & e_0^{(n_2)}
\end{pmatrix}, \quad H_{d} =
\begin{pmatrix} H_{\text{BOTTOM}} & e_0^{(n_2)} & Z_{n_2} r_b
\end{pmatrix} $$

\vspace{1em}
\subsubsection{Packing}

The pack operation reverts the changes of the splitting phase.
After computing the inverse sub-blocks (denoted as $x, y, z, t$),
our goal is to assemble the global generators $G_{INV}$ and $H_{INV}$
for the full inverse matrix $A^{-1}$.
If we were to concatenate directly the local generators,
the result would assume that the boundaries between the sub-blocks
are filled with zeros.
Because the global operator $\Phi_+$ shifts elements,
we must correct the bleeding to restore data integrity.
Just as we calculated boundaries during the split, we must extract
the boundary vectors of the local inverse blocks.
Since we are now operating on isolated sub-block generators, the
coordinates for each reconstruction are local and starts from zero.

To apply these corrections, we construct a large set of global generators.
We horizontally concatenate the local generators of $x, y, z, t$ into
their respective quadrants and append four specific boundary
correction columns ($c_1, c_2, c_3, c_4$).
We define these global assembled generators as block matrices:
$$ G_{RES} =
\begin{pmatrix}
  G_x & G_y & 0 & 0 & -Z_{n_1} v_x & 0 & 0 & 0 \\
  0 & 0 & G_z & G_t & 0 & -e_0^{(n_2)} & -(s_x e_0^{(n_2)} + Z_{n_2}
  v_z) & -e_0^{(n_2)}
\end{pmatrix} $$
$$ H_{RES} =
\begin{pmatrix}
  H_x & 0 & H_z & 0 & 0 & Z_{n_1} r_x & 0 & 0 \\
  0 & H_y & 0 & H_t & e_0^{(n_2)} & 0 & e_0^{(n_2)} & Z_{n_2} r_y
\end{pmatrix} $$

Notice that the boundary correction columns in $G_{RES}$ carry negative signs.
The reconstruction is going to add the values diagonally, to correct
bleeding, we must
substract the the corrections.
To prove the correctness of these columns,
we compute the global displacement reconstruction by multiplying
$\Phi_+(A^{-1}) = G_{RES} H_{RES}^T$:

$$ G_{RES} H_{RES}^T =
\begin{pmatrix}
  G_x H_x^T                               & G_y H_y^T - Z_{n_1} v_x
  (e_0^{(n_2)})^T \\
  G_z H_z^T - e_0^{(n_2)} (Z_{n_1} r_x)^T & G_t H_t^T - Z_{n_2} v_z
  (e_0^{(n_2)})^T - e_0^{(n_2)} (Z_{n_2} r_y)^T - s_x e_0^{(n_2)}
  (e_0^{(n_2)})^T
\end{pmatrix} $$

The isolated terms ($G_x H_x^T, G_y H_y^T$, etc.) represent the
isolated local displacements of the sub-matrices.
The subtracted terms are the bleeds.
For example, the term $- Z_{n_1} v_x (e_0^{(n_2)})^T$ subtracts the
shifted right column of block $x$ as it crosses into the first column
of block $y$.

This structural reassembly exposes a vulnerability. Let $\alpha_x,
\alpha_y, \alpha_z, \alpha_t$ be the displacement ranks of the sub-quadrants.
The rank of our assembled global generator is:
$$ \text{Rank}_{RES} = \alpha_x + \alpha_y + \alpha_z + \alpha_t + 4 $$
Passing this uncompressed matrix up to the next recursion level would
cause the displacement rank to explode,
ruining the optimisations we aimed for at the first place. We call
the compression seen in section \ref{ch:compression} to keep the
algorithm efficient.

\subsection{Inversion Benchmarks}

The performance benchmarks for the generator-based Strassen inversion
show improvements over standard dense operations at larger scales:

\begin{table}[H]
  \centering
  \begin{tabular}{r r r}
    \hline
    $n$  & Generators (ms) & Flint Dense (ms) \\
    \hline
    128  & 5.328           & 1.967 \\
    256  & 12.79           & 12.03 \\
    512  & 30.74           & 80.16 \\
    1024 & 74.50           & 554.8 \\
    2048 & 180.4           & 3975 \\
    4096 & 452.9           & 28300 \\
    \hline
  \end{tabular}
  \caption{Average runtimes of inversion operations ($\alpha =
    2$). Averages are computed over 100 iterations for $n \le 512$, 20
  iterations for $n \in \{1024, 2048\}$, and 10 iterations for $n = 4096$.}
  \label{tab:bench_inv_size}
\end{table}

\begin{table}[H]
  \centering
  \begin{tabular}{l r r r}
    \hline
    Method | $\alpha$        & $n=512$ (ms) & $n=1024$ (ms) & $n=2048$ (ms) \\
    \hline
    Generators ($\alpha=2$)  & 34.65        & 77.96         & 180.4 \\
    Generators ($\alpha=4$)  & 61.12        & 144.8         & 340.2 \\
    Generators ($\alpha=8$)  & 135.5        & 333.5         & 800.1 \\
    Generators ($\alpha=16$) & 363.8        & 931.5         & 2259.0 \\
    Generators ($\alpha=32$) & 1064.0       & 2829.0        & 7313.0 \\
    \hline
    Flint Dense (Baseline)   & 80.16        & 554.8         & 3975.0 \\
    \hline
  \end{tabular}
  \caption{Average runtimes of generator inversion scaling by rank
  over 10 runs, compared to constant dense inversion baseline.}
  \label{tab:bench_inv_rank}
\end{table}

We analyze the execution time as $n$ grows.
While the split and pack
operate in $O(\alpha M(n))$. The recursive
algorithm involves multiplying matrices with generators seen in
section \ref{section:matrixmatrixmul} and thus
overall algorithm is bounded by $O(\alpha^2 M(n) \log n)$.

Taking $M(n) = O(n \log n)$, the theoretical
complexity is $O(\alpha^2 n \log^2 n)$. For a
constant rank, the expected scaling factors when increasing $n$ from
128 to 1024 ($8\times$) and from 1024 to 4096 ($4\times$) are
approximately 16.32 and 5.76 respectively. The
observed scaling ratios are:

$$ \frac{T(1024)}{T(128)} \approx 14.43, \quad
\frac{T(4096)}{T(1024)} \approx 6.02 $$

These observed ratios seem to align with the
expectations.
The slight deviation at smaller sizes (14.43 and 16.32)
can be explained by the recursive overhead and the constant factors
of copying matrix elements into polynomials. Also the difference at larger
sizes (6.02 and 5.76) can be explained as the data structure
scales up, they begin to exceed the CPU's fast cache limits.
Eventhough, as $n$ progresses, complexity respects the
$O(\alpha^2 M(n) \log n)$ bound.
When comparing against the dense approach, FLINT follows an
$O(n^3)$ curve to over 28 seconds at $n=4096$.
Inversion with generators completes the same operation in just over
450 milliseconds, giving us a $\sim \! 62\times$ speedup for a
Toeplitz matrix.
But this speedup will decrease as rank increases on a quasi-Toeplitz
case, as the efficiency of storing a dense matrix decreases too.
However, the results reveal an intersection point. For
small matrices ($n \le 256$), the recursive overhead and cost of
splitting and packing the sub matrices make the algorithm
slower than dense inversion. As seen in table
\ref{tab:bench_inv_size}, the crossover occurs between $n=256$ and $n=512$.

Table \ref{tab:bench_inv_rank} highlights the algorithm's dependency
to the displacement rank $\alpha$. Since the overall complexity
contains linear operations (splitting and
packing in $O(\alpha)$) and quadratic operations (generator
multiplication in $O(\alpha^2)$), we expect the $\alpha^2$ term as
$\alpha$ grows.
Doubling the rank from $2$ to $4$ increased the runtime by roughly
$1.85\times$, but doubling it from $16$ to $32$ increased the runtime
by a factor of over $3.2\times$. Because the algorithm scales quadratically
any doubling of the rank forces the algorithm to do 4 times as much work
This confirms that to maintain high
performance in iterative applications, periodic generator compression
(as detailed in Section \ref{ch:compression}) is absolutely critical.

\section{Conclusion and Perspectives}

In this project, we successfully implemented
arithmetic algorithms for quasi-Toeplitz matrices using displacement
generators. However, while Strassen-type inversion is highly
effective for large matrices with small ranks, the rank explosion
problem and computational overhead limit its viability for smaller
systems or highly iterative applications.

For future optimizations, alternative algorithmic structures could be
preferred. As detailed by Victor Y. Pan \cite[Chapter
5]{pan2001structured}, the Morf-Bitmead-Anderson (MBA) algorithm
achieves an even faster $O(n \log^2 n)$ complexity for
Toeplitz-like matrices. However, as it is also a divide-and-conquer
method, it suffers from similar calculations and
generator compression. To bypass matrix splitting
entirely, the Newton iterative version based on the Schulz iteration
$X_{k+1} = X_k(2I - A X_k)$—is highly preferable \cite[Chapters 6 and
7]{pan2001structured}. If the objective is to solve a linear system
$Ax = b$, computing the
full inverse $A^{-1}$ introduces unnecessary computations.
Instead, one can take direct advantage of the structured properties
of $A$ to solve the system more efficiently.
Methods such as the Levinson-Durbin algorithm
\cite{durbin1960fitting}, avoid inversion and solve Toeplitz systems
in $O(n^2)$ operations.
Iterative Krylov subspace methods, the Preconditioned Conjugate
Gradient (PCG) method \cite{doi:10.1137/0909051} use FFT to compute
matrix-vector products, achieving $O(n \log n)$ complexity per iteration.
These iterative methods they rely on numerical approximation and
convergence tolerances.
For an exact deterministic algebraic solution, there is an
alternative by resolving via M-Padé
approximations\cite{khichane2026matricesdisplacementstructuredeterministic},
which achieves a deterministic $\tilde{O}(\alpha^{\omega-1}n)$ complexity.
